% V18 integration of Dot batch 088; immutable source SHA256 02f14bc6905edd6df5b519ec0a61968ce9079b424b003c5ffac78dd853d96a08
\begingroup
% These definitions have only group-local scope; the shared preamble is untouched.
\def\OO{\mathcal O}\def\FF{\mathcal F}\def\GG{\mathcal G}
\def\PP{\mathbf P}\def\calP{\mathcal P}\def\dd{\delta}
\def\Div{\operatorname{Div}}\def\depth{\operatorname{depth}}
\def\pr{\operatorname{pr}}\def\coker{\operatorname{coker}}
\par\medskip\noindent\textbf{原手稿 IV.20：超平面截面与锥形投影。}
\egasection{§1\quad 预备知识与记号（法文原号20.1）}

设 \(S\) 为预概形，\(E\) 为 \(S\) 上有限型局部自由模，\(E^\vee\) 为其对偶。
记 \(P=\PP(E)\) 为 \(E\) 所定义的射影丛，\(P^\vee=\PP(E^\vee)\)。
我们称 \(P^\vee\) 为 \(P\) 的\textbf{超平面概形}。以下说明这个名称。\footnote{法物理3页章题原打字词“线性”已划去，手改为“锥形”。两个配对式的法原标签为单叉、双叉，英编稿为星号、双星号，以下沿用既有译本的星号。}

给定 \(P^\vee\) 在 \(S\) 上的截面 \(\xi\)，便有 \(E^\vee\) 的一个可逆商模 \(L\)。
在 \(P\) 上，既有 \((E_P)^\vee\) 的可逆商 \(L_P\)，也有 \(E_P\) 的可逆商
\(\OO_P(1)\)。取对偶后，可将 \(L_P^{-1}\) 与 \(\OO_P(-1)\) 分别视作
\(E_P\) 与 \((E_P)^\vee\) 的可逆子模，且局部为直和因子。
评价配对 \(E_P\otimes(E_P)^\vee\to\OO_P\) 给出
\begin{equation}
 \OO_P(-1)\otimes L_P^{-1}\longrightarrow\OO_P,\tag{*}
\end{equation}
或其转置
\begin{equation}
 \OO_P\longrightarrow\OO_P(1)\otimes L_P=L_P(1).\tag{**}
\end{equation}
这就是由 \(\xi\) 典范确定的 \(L_P(1)\) 的截面。
这个截面的“除子”，即由 (*) 的像理想定义的闭子概形 \(H_\xi\subset P\)，
称为 \(\xi\in P^\vee(S)\) 定义的\textbf{超平面}。

在 \(S\) 上局部地，\(\xi\) 可由 \(E\) 的一个在每个点的纤维中均不为零的截面
\(\phi\) 给出；改变代表只会乘上 \(\OO_S\) 的一个可逆截面。
对投影 \(p:P\to S\)，有 \(E=p_*\OO_P(1)\)，
因而 \(\phi\) 也给出 \(\OO_P(1)\) 的截面，其除子正是 \(H_\xi\)。
把 \(L^{-1}\) 视为 \(E\) 的局部直和因子时，取它的处处不为零的截面（也就是一个局部平凡化）
便给出这种 \(\phi\)。又有典范同构
\[
 H_\xi\simeq \PP(E/L^{-1})\subset \PP(E).
\]
这是第三种描述。它也表明 \(H_\xi\) 是 \(S\) 上的射影丛，因而在 \(S\) 上光滑。
（作者提示：应当已在第17段说明射影丛的光滑性。）或许从这种描述出发更好。
这里按射影丛的约定把秩零商的射影化看成空概形。\footnote{例如 \(E\) 秩一时，\(E/L^{-1}=0\)，对应的超平面为空，仍是光滑的。下文以次数区分相对除子分支时，不能沿用这一退化情形下没有区别的“次数”。}

\noindent\textbf{备注。}
\(H_\xi\) 的形成与基变换相容，给出集合值反变函子的自然变换
\[
 P^\vee\longrightarrow \Div(P/S).
\]
右端为 \(P/S\) 的相对除子函子：在任意 \(S\)-预概形 \(S'\) 上，
其值是 \(P_{S'}\) 中相对于 \(S'\) 为余维一的横截完全交的闭子概形集合（参看第19段）。
此变换为单态射，即 \(\xi\) 由 \(H_\xi\) 确定。
作者在此预告相对除子函子的表示。按商线约定订正其参数模后，在 $E$ 的秩处处至少为二的情形，
相应的各次数参数丛应写为 $\PP((\operatorname{Sym}^k E)^\vee)$；作者预告它们的不交并表示该函子，
于是 \(P^\vee\) 就是其中既开又闭的次数一部分。
（作者记：相对除子的确定现在就可借用\textbf{第III章}的结果，
作为例子补到第19段。）\footnote{法物理4页清楚写的是对称幂的上标 \(k\)，没有给出以域 \(k\) 为底的张量记号；本段的底始终是 \(S\)。同页末行写第III章，英编稿作第II章，旧中文从英稿，现据法稿订正。法稿原式为 \(\PP(\operatorname{Sym}^k(E^\vee))\)，对偶放在 $E$ 上。按 $p_*\OO_P(k)=\operatorname{Sym}^k E$ 的商线约定，次数 $k$ 的方程给出 $\operatorname{Sym}^k E$ 中的一条局部直和因子线，故其参数模是 $(\operatorname{Sym}^k E)^\vee$。任意底上不能无条件把对偶与对称幂交换；$k!$ 可逆时的对称化识别不能直接用于任意特征。正文只订正这一预告的参数模，$k=1$ 的超平面构造不受影响。秩一的射影丛不能把多个同样的空除子按各个 \(k\) 重复表示。这里未证明一般相对除子函子的表示性，旧注 P03 的历史记录仍保留在核查报告中。}

作基变换 \(S'=P^\vee\to S\)，取 \(P^\vee_{S'}\) 在 \(S'\) 上的对角截面（“泛截面”），
得到 \(P\times_S P^\vee\) 中的闭子概形 \(H\)，称为 \(P\) 与 \(P^\vee\) 的\textbf{关联概形}，
其理想为典范同态
\[
 \OO_P(-1)\boxtimes_S\OO_{P^\vee}(-1)
 \longrightarrow \OO_{P\times_S P^\vee}
\]
的像。它在 \(P^\vee\) 上是射影丛；由对称性，在 \(P\) 上也是射影丛。
沿 \(\xi:S\to P^\vee\) 对“普遍超平面” \(H\) 作基变换，就恢复特殊超平面 \(H_\xi\)。
同样，可对任意 \(S'\)-值点 \(\xi:S'\to P^\vee\) 定义 \(H_\xi\subset P_{S'}\)。

设 \(X\) 是在 \(S\) 上有限型的预概形，\(f:X\to P\) 是 \(S\)-态射。\footnote{法物理5页此行右端在“sur”后出框，英稿正文写“在 \(P\) 上有限型”并附编者疑问。两种读法在这里等价：\(P\to S\) 有限型且分离；若 \(X\to S\) 有限型，则 \(f\) 由闭图嵌入 \(X\to X\times_S P\) 和对 \(P\) 的有限型投影分解；反向由复合即得。}
本段的一项主要目标是对每个超平面 \(H_\xi\subset P\) 研究它的逆像
\[
 Y_\xi=f^{-1}(H_\xi)=X\times_P H_\xi,
\]
特别是把 \(X\) 与 \(Y_\xi\) 的性质联系起来。和通常一样，也要考虑任意 \(S'\)-值点：
\[
 Y_\xi=f_{S'}^{-1}(H_\xi)=X_{S'}\times_{P_{S'}}H_\xi
       =X\times_P H_\xi .
\]
最后一项是把 \(H_\xi\) 通过 \(H_\xi\to P_{S'}\to P\) 当成 \(P\)-预概形。
取 \(S'=P^\vee\) 与对角截面，记 \(Y=X\times_P H\)
（作者括注：除非迪厄多内建议更好的记号）。
一般地，\(Y_\xi=Y\times_{P^\vee}S'\)。
设 \(\FF\) 为 \(X\) 上的模层，\(\GG\) 为它在 \(Y\) 上的拉回，
\(\GG_\xi\) 为它在 \(Y_\xi\) 上的拉回。若 \(q_\xi:Y_\xi\to Y\) 是投影，
则 \(\GG_\xi\simeq q_\xi^*\GG\)。这是原稿所写张量式
\(\GG_\xi=\GG\otimes_{\OO_{P^\vee}}\OO_{S'}\) 按基变换读出的含义。\footnote{法物理6页及英6正文都把 \(\GG\) 的所在写成 \(H\)，同页原图却把 \(\GG\) 写在 \(Y\) 上；一般也不存在所需的 \(H\to X\)。据原图与拉回复合改为 \(Y\)。英译编注关于模层是否拟凝聚或凝聚的疑问保留为编注，不在此额外加假设。}

主要构造如下，\(\FF,\GG,\GG_\xi\) 分别在 \(X,Y,Y_\xi\) 上：
\[
\begin{tikzcd}[column sep=1.5em,row sep=2em]
X\arrow[d,"f"'] & X\times_S P^\vee\arrow[l]\arrow[d]
   &Y\arrow[l,hook']\arrow[d]&Y_\xi\arrow[l]\arrow[d]\\
P\arrow[d] & P\times_S P^\vee\arrow[l]\arrow[d]
   &H\arrow[l,hook']\arrow[dl]&H_\xi\arrow[l]\arrow[dl]\\
S & P^\vee\arrow[l]& S'\arrow[l,"\xi"']&
\end{tikzcd}
\]
图中的方形与菱形均为笛卡尔图。
下一节研究 \(S'=\Spec K\) 为域的谱，其像为相应纤维 \(P_s^\vee\) 的泛点。
沿 \(\Spec k(s)\to S\) 基变换后，可令 \(S=\Spec k\)，其中 \(k\) 是域。
研究的大多数性质是几何性质，可视所考虑的几何问题作域扩张。
可采用 \(K=k(\eta)\) 与到 \(P^\vee\) 的典范态射，其中 \(\eta\) 是泛点，
或在适当基变换后使用代数闭域。对 \(X,Y_\xi\) 的几何问题，
也可据相应的基变换把 \(k\) 换成代数闭域。
若 \(f\) 是浸入，通常称 \(Y_\xi\) 为 \(X\) 的\textbf{超平面截面}
（相对于 \(f\) 和 \(H_\xi\)）；这一术语也可以用于任意 \(f\)。

\egasection{§2\quad 泛超平面截面：局部性质（法文原号20.2）}

现令 \(S=\Spec k\)，\(k\) 为域。
对 \(P^\vee\) 的点 \(\eta\) 及典范态射 \(\xi:\Spec k(\eta)\to P^\vee\)，
也用 \(H_\eta,Y_\eta,\GG_\eta\) 分别代替 \(H_\xi,Y_\xi,\GG_\xi\)。
本节后面 \(\eta\) 始终指 \(P^\vee\) 的泛点。\footnote{法物理7页在 \(\Spec\) 后直接写点符号，英稿展开为 \(\Spec k(\eta)\)。这里采用剩余域的明确写法。命题2.1旁的手写补句可辨为“此外 \(Y\) 不可约”，英编稿夹有“字迹不清，问 AG”及补入文字；该补句已由下述证明直接给出。}

\egas{2.1}{命题}
若 \(X\) 不可约，则 \(Y_\eta\) 不可约或为空；非空时它占优于 \(X\)。
\begin{proof}
因 \(H\to P\) 是射影丛，\(Y\to X\) 也是。若 \(Y\) 非空，
它是不可约 \(X\) 上的非空射影丛，故不可约并占优于 \(X\)。
当 \(Y_\eta\ne\varnothing\) 时，它是不可约 \(Y\) 在不可约 \(P^\vee\) 上的非空泛纤维。
其泛点正是 \(Y\) 的泛点，位于 \(X\) 的泛点上方，故结论成立。
\end{proof}

\egas{2.2}{命题}
设 \(Z\subset P\) 为点集。其在 \(H_\eta\) 中的逆像 \(Z_\eta\) 为空，
当且仅当 \(Z\) 的每个点在 \(P\) 中均闭。
特别地，若 \(Z\) 可构造，则 \(Z_\eta=\varnothing\) 当且仅当 \(Z\) 有限。
\begin{proof}[原稿论证与译校展开]
先取 \(Z=\{z\}\)，令 \(X=\overline{\{z\}}\subset P\) 取约化闭子概形结构。
由（2.1），若 \(X\) 的泛超平面截面非空，便有点投到 \(z\)；
只需证明该截面为空当且仅当 \(X\) 在 \(k\) 上有限。
若截面为空，\(X_{k(\eta)}\) 既为射影概形，又是超平面补集这一仿射概形的闭子概形。
仿射且固有的有限型态射为有限态射；于是 \(X_{k(\eta)}\) 有限，有限性沿域扩张下降。
反之，若 \(X\) 有限，记 \(n=\dim P\)。在 \(n\ge1\) 时，
\(Y\to X\) 的相对维数为 \(n-1\)，故 \(Y\) 的维数至多为 \(n-1\)；
不能占优于维数 \(n\) 的 \(P^\vee\)；若 \(n=0\)，关联概形为空。
最后，有限型 \(k\)-概形中由闭点组成的可构造集必有限：
在其闭包的每个不可约分支中，该集包含非空相对开集，因而包含分支泛点；
该泛点既是闭点，分支只能为一点。有限可构造集中的所有点亦闭，
因为有限点集的任一点可用有限个开集及其补集同其余点区分，故各单点可构造；单点可构造即局部闭，而有限型 \(k\)-概形中局部闭点皆闭。
\end{proof}
\noindent\emph{译注。} 法物理7页证明中“\(H\) 在 \(P\) 中的像”须取 \(H_\eta\)，
否则 \(n\ge1\) 时 \(H\to P\) 满射。本稿使用了泛纤维，并展开原稿括注的“仿射且射影”论证。

\egas{2.3}{推论}
设 \(f:X\to P\) 有限型，\(Z\subset X\) 可构造。
其在 \(Y_\eta\) 中的逆像为空，当且仅当 \(f(Z)\) 有限。
特别地，\(Y_\eta=\varnothing\) 当且仅当 \(f(X)\) 有限。
这是将（2.2）用于可构造像 \(f(Z)\)；纤维积上剩余域的张量积非零，
故只要 \(x\in Z\) 与 \(h\in H_\eta\) 在 \(P\) 有同像，就有上方的点。

\egas{2.4}{推论}
设 \(Z,Z'\) 是 \(X\) 的两个闭子集，且 \(Z\) 不可约。
分别以 \(Z_\eta,Z'_\eta\) 表示它们在 \(Y_\eta\) 中的逆像。则
\[
 Z_\eta\subseteq Z'_\eta
 \quad\Longleftrightarrow\quad
 Z\subseteq Z'\ \text{或}\ f(Z)\ \text{有限}.
\]
而在不另要求 \(Z'\) 不可约时，正确的等号判据是
\[
 Z_\eta=Z'_\eta\quad\Longleftrightarrow\quad
 f(Z\mathbin{\triangle} Z')\ \text{有限}.
\]
这里 \(\triangle\) 为对称差。若另外假设 \(Z'\) 也不可约，就可把等号判据改写为
“\(Z=Z'\)，或 \(f(Z),f(Z')\) 均有限”。\footnote{法物理8页原打字“不可约”修饰两个闭集，已划去并手补“\(Z\) 不可约”；原等号句却没有同步修改：它仍称 $Z_\eta=Z'_\eta$ 当且仅当 $Z=Z'$ 或 $f(Z),f(Z')$ 都有限。取 \(X=P=\PP^2_k\)、\(f=\mathrm{id}\)、\(Z\) 为一条直线，取线外一个 \(k\)-有理点 \(x\)，令 \(Z'=Z\cup\{x\}\)。则 \(Z_\eta=Z'_\eta\)，但 \(Z\ne Z'\)，且两像都无限。故不能按稿上修改后的假设保留原等号句。}
\begin{proof}[原包含判据的证明及译校补证]
逆像与集合差及对称差相容，由（2.3）得到等号判据，并知
\(Z_\eta\subseteq Z'_\eta\) 等价于 \(f(Z\setminus Z')\) 有限。
若 \(Z\nsubseteq Z'\)，则 \(Z\setminus Z'\) 稠密于 \(Z\)，所以它的像在 \(f(Z)\) 中稠密。
该像若有限，因又可构造，便是 \(P\) 中的有限闭点集，于是 \(f(Z)\) 亦有限。
反向是直接代入（2.3）。两边均不可约时，对两种包含分别应用第一判据：
若 \(Z\ne Z'\) 且 \(Z\subsetneq Z'\)，由反向包含得 \(f(Z')\) 有限，从而 \(f(Z)\) 有限；
真反包含情形类似；互不包含时两像各自有限。
\end{proof}

\egas{2.5}{推论}
对每个满足 \(\dim\overline{f(X_i)}>0\) 的不可约分支 \(X_i\) 取其在 \(Y_\eta\) 中的逆像
\(Y_{i\eta}\)。这样得到这些 \(X_i\) 与 \(Y_\eta\) 的不可约分支之间的一一对应。
\begin{proof}
由（2.3）及有限型假设，\(Y_\eta\) 是这些闭而非空的逆像的有限并。
由（2.1）这些逆像不可约，由（2.4）的\emph{包含}判据两两互不包含。
所以它们恰为 \(Y_\eta\) 的不可约分支。
\end{proof}

注意，若 $\dim X_i=d_i$，则对上述非空分支有 $\dim Y_{i\eta}=d_i-1$。更一般地：\cmt{D18-088-P25}{这是法文物理第8页（2.5）证明之后的过渡句；原稿写 $Y_i$，依泛截面语境补明下标 $\eta$。第88批将其解释并入（2.6）脚注，本次恢复其独立位置。}

\egas{2.6}{命题}
若每个不可约分支 \(X_i\) 均满足 \(\dim\overline{f(X_i)}>0\)，
即 \(Y_{i\eta}\ne\varnothing\)，则
\[
 \dim Y_\eta=\dim X-1.
\]
若 \(f\) 是浸入且 \(\dim\overline{f(X)}>0\)，同样有这一等式。\footnote{法物理8页边批另写浸入情形。英编稿对连词有疑问。此处以两个经核对足够的假设分写；浸入情形的证明无需误说零维分支的泛超平面截面非空。既有译文中的“\(\dim Y_i=d_i-1\)”应当指 \(Y_{i\eta}\)，不是关联族总空间。}
\begin{proof}[原稿论证与译校展开]
对每个有非空截面的分支，只须在其约化结构下计算维数。设该分支为不可约约化的 \(X_i\)。
由于 \(k(\eta)/k\) 纯超越，\(X_{i,k(\eta)}\) 仍为整概形。
\(Y_{i\eta}\) 由 \(f^*\OO_P(1)\) 的基变换后截面切出。
该截面不恒零，因为 \(Y_{i\eta}\) 在 \(X_i\) 的像避开闭点（用（2.3）），
而 \(X_{i,k(\eta)}\to X_i\) 忠实平坦。故在这个整概形上切出非空有效卡蒂埃除子，
各不可约分支均余维一。有限型域上维数公式与域扩张不改变维数给出
\[
 \dim Y_{i\eta}=\dim X_{i,k(\eta)}-1=\dim X_i-1.
\]
最后取非空截面分支维数的最大值。在浸入情形，最大维数分支为正维，
其像闭包维数与分支维数相同，故它的截面不会空；
即使存在零维分支，其空截面也不影响最大值。
\end{proof}
\noindent\emph{译注。} 法物理9页在维数公式处留有参考待补；
这里把约化、非零截面和取最大值的步骤写清，不称原稿已给出所引文献。

\egas{2.7}{命题}
设 \(\FF\) 为 \(X\) 上的拟凝聚模。取它的关联素循环中所有满足
\(\dim\overline{f(Z_i)}>0\) 的 \(Z_i\)，令 \(Z_{i\eta}\) 为其在 \(Y_\eta\) 中的逆像。
它们恰为 \(\GG_\eta\) 的全部关联素循环，包含关系与对应 \(Z_i\) 的包含关系相同。
\begin{proof}
包含关系由（2.4）的包含判据得出。
\(Y\to X\) 是射影丛，因而平坦，非空纤维不可约且满足 \((S_1)\)。
原稿援引第3段的平坦拉回结论：\(\GG\) 的关联素循环由 \(\FF\) 的关联素循环取逆像得到。
再限制到泛纤维 \(Y_\eta\)，只保留非空的逆像；
由（2.3），它们正对应所列的 \(Z_i\)。
原稿并指出，可直接利用 \(Y_\eta\to X\) 平坦、非空纤维几何正则且几何不可约：
它们来自 \(Y\to X\) 的射影空间纤维的局部化。这个观察也可用于（2.1）。
\end{proof}
\noindent\emph{译注。} “局部化”指为取 \(P^\vee\) 的泛纤维而沿底的泛点局部化，
不是把任意闭纤维都认成整个射影空间。此处保留作者所用的关联素循环术语及第3段引用；
没有把版本未标明的原引文改成新定理编号。

\endgroup
